\documentclass[12pt,a4paper]{article}

% -------------------------------------------------
% Pakete
% -------------------------------------------------
\usepackage[utf8]{inputenc}
\usepackage[T1]{fontenc}
\usepackage[ngerman]{babel}
\usepackage[a4paper,margin=2.5cm]{geometry}
\usepackage{booktabs}
\usepackage{tabularx}
\usepackage{longtable}
\usepackage{xcolor}
\usepackage{listings}
\usepackage{hyperref}
\usepackage{enumitem}
\usepackage{parskip}

\hypersetup{
    colorlinks=true,
    linkcolor=blue,
    urlcolor=blue
}

\definecolor{codegray}{RGB}{245,245,245}

\lstset{
    backgroundcolor=\color{codegray},
    basicstyle=\ttfamily\small,
    breaklines=true,
    frame=single,
    columns=fullflexible,
    showstringspaces=false
}

% -------------------------------------------------
% Dokument
% -------------------------------------------------
\begin{document}

\begin{titlepage}
    \centering

    \vspace*{3cm}

    {\Huge\textbf{Integration von Twilio für WhatsApp und SMS}}\\[0.8cm]

    {\Large Technisches Konzept für den Nachrichtenversand}\\[2cm]

    \rule{13cm}{0.5pt}\\[0.5cm]

    {\large Projekt: Workflowbasierte Rückbestätigung von Lieferterminen}\\[0.3cm]
    {\large Autorin: Buse Okcu}\\[0.3cm]
    {\large Datum: \today}\\[0.5cm]

    \rule{13cm}{0.5pt}

    \vfill
\end{titlepage}

\tableofcontents
\newpage

% =================================================
\section{Zielsetzung}
% =================================================

Im Projekt besteht bereits ein E-Mail-Versand, der in der lokalen
Entwicklungsumgebung über Mailpit getestet werden kann. Aufbauend auf
dieser bestehenden Kommunikationsfunktion soll das System um zwei weitere Kanäle erweitert werden: SMS und WhatsApp.

Ziel ist es, Kundenbenachrichtigungen und
Liefertermin-Bestätigungsanfragen nicht nur per E-Mail, sondern
zusätzlich auch über mobile Nachrichtenkanäle bereitzustellen.

Diese Dokumentation beschreibt die geplante technische Erweiterung des
bestehenden Kommunikationsmodells um \textbf{Twilio} als zentralen
Provider für SMS- und WhatsApp-Nachrichten. Dabei werden die relevanten
Felder, Konfigurationswerte und die vorgesehene Nachrichtenstruktur
erläutert.

% =================================================
\section{Warum Twilio}
% =================================================

Im Rahmen der Technologieauswahl wurden auch alternative
Messaging-Lösungen betrachtet und mit Twilio verglichen. Ausschlaggebend
für Twilio war insbesondere, dass bereits in der Testphase mit einem
kostenlosen Trial-Setup beziehungsweise einer Twilio-Testnummer erste
Nachrichtenflüsse erprobt werden können. Zusätzlich stellt Twilio mit
dem \emph{Content Template Builder} eine Benutzeroberfläche bereit, über
die Nachrichtenvorlagen für den Versand verwaltet und erstellt werden
können. Dadurch lassen sich wiederverwendbare Templates nicht nur per
API, sondern auch komfortabel über die Konsole pflegen.

Die Wahl von Twilio bietet folgende Vorteile:

\begin{itemize}
    \item Gemeinsame API für WhatsApp und SMS
    \item Unterstützung von WhatsApp-Sandbox für Entwicklungszwecke
    \item Versand von SMS über dieselbe Plattform
    \item Frühe Testmöglichkeiten über ein kostenloses Trial-Setup
    \item Erstellung und Verwaltung von Nachrichtentemplates über die Twilio-Oberfläche
\end{itemize}

Dadurch kann die Anwendung beide Kommunikationskanäle mit einer
ähnlichen Integrationslogik ansprechen, ohne für jeden Kanal einen
vollständig separaten Dienst einführen zu müssen.

% =================================================
\section{Einsatzszenario}
% =================================================

Im Projekt wird Twilio als externer Messaging-Provider eingesetzt.
Das Backend erzeugt aus den vorliegenden Auftrags- und Kundendaten eine
Nachricht und übermittelt diese an die Twilio Messaging API. Twilio
übernimmt anschließend die Zustellung an die angegebene Zielrufnummer.

Dadurch wird die bereits vorhandene E-Mail-Kommunikation um die
Kommunikationskanäle SMS und WhatsApp erweitert, ohne dass die
bestehende Fachlogik grundlegend angepasst werden muss.

% =================================================
\section{Technische Architektur}
% =================================================

Die Anwendung verwaltet ihre Geschäftslogik intern unabhängig von
Twilio. Für den Versand wird jedoch ein Adapter vorgesehen, der die
internen Daten in Twilio-kompatible Requests übersetzt.

Der geplante Ablauf ist wie folgt:

\begin{enumerate}
    \item Das Backend erhält einen neuen Kommunikationsauftrag.
    \item Der gewünschte Kanal wird ermittelt, zum Beispiel SMS oder
    WhatsApp.
    \item Aus den vorhandenen Auftrags- und Kundendaten wird eine
    Nachricht erzeugt.
    \item Die Nachricht wird an die Twilio API gesendet.
    \item Twilio übernimmt die Zustellung an die Zielrufnummer.
\end{enumerate}

Dadurch bleibt die Fachlogik von der externen API weitgehend getrennt
und kann später bei Bedarf auf andere Anbieter erweitert werden.

% =================================================
\section{Benötigte Konfigurationswerte}
% =================================================

Für die Anbindung an Twilio werden insbesondere folgende
Konfigurationswerte benötigt:

\begin{table}[h]
    \centering
    \begin{tabularx}{\textwidth}{l l X}
        \toprule
        \textbf{Konfiguration} & \textbf{Typ} & \textbf{Beschreibung} \\
        \midrule
        TWILIO\_ACCOUNT\_SID
        & String
        & Eindeutige Konto-ID des Twilio-Accounts \\

        TWILIO\_AUTH\_TOKEN
        & String
        & Authentifizierungs-Token für API-Aufrufe \\

        TWILIO\_SMS\_FROM
        & String
        & Twilio-Telefonnummer für den SMS-Versand \\

        TWILIO\_WHATSAPP\_FROM
        & String
        & WhatsApp-Absenderkennung, z.\,B. \texttt{whatsapp:+14155238886} in der Sandbox \\

        TWILIO\_CONTENT\_TEMPLATE\_SID
        & String
        & SID eines gemeinsamen Twilio-Content-Templates für SMS und WhatsApp \\
        \bottomrule
    \end{tabularx}
    \caption{Erforderliche Konfigurationswerte}
\end{table}

% =================================================
\section{Relevante Fachfelder im Projekt}
% =================================================

Für den Versand einer Bestätigungsanfrage werden im Projekt bereits
fachliche Daten übergeben, aus denen anschließend die jeweilige
E-Mail-, SMS- oder WhatsApp-Nachricht erzeugt werden kann.

Die relevanten Felder sind dabei:

\begin{table}[h]
    \centering
    \begin{tabularx}{\textwidth}{l l X}
        \toprule
        \textbf{Feld} & \textbf{Datentyp} & \textbf{Beschreibung} \\
        \midrule
        externalOrderId
        & String
        & Eindeutige externe Auftragsnummer \\

        customerName
        & String
        & Name der Kundin oder des Kunden \\

        communicationChannel
        & Enum
        & Gewählter Kommunikationskanal, z.\,B. EMAIL, SMS oder WHATSAPP \\

        customerEmail
        & String
        & E-Mail-Adresse der Kundin oder des Kunden \\

        customerPhoneNumber
        & String
        & Telefonnummer im E.164-Format für SMS oder WhatsApp \\

        deliveryAddress
        & String
        & Lieferadresse \\

        product
        & String
        & Bestelltes Produkt \\

        quantityLiters
        & Integer
        & Bestellte Menge in Litern \\

        deliveryDate
        & LocalDate
        & Geplantes Lieferdatum \\

        deliveryWindowStart
        & LocalTime
        & Beginn des Lieferzeitfensters \\

        deliveryWindowEnd
        & LocalTime
        & Ende des Lieferzeitfensters \\

        responseDeadlineHours
        & Integer
        & Frist in Stunden für die Rückmeldung \\

        priceDisplayText
        & String
        & Preisangabe zur Anzeige für die Kundin oder den Kunden \\
        \bottomrule
    \end{tabularx}
    \caption{Fachliche Eingabefelder für eine Kommunikationsanfrage}
\end{table}

% =================================================
\section{Twilio-Requests}
% =================================================

\subsection{Gemeinsames Nachrichtenmodell für SMS und WhatsApp}

Sowohl für SMS als auch für WhatsApp wird im Projekt die
Twilio Messages API verwendet. Beide Kanäle nutzen dabei
dieselbe Template-Struktur über ein gemeinsames
Twilio Content Template.

Der technische Aufbau der Requests ist daher weitgehend identisch.
Unterschiede ergeben sich im Wesentlichen nur beim jeweiligen
Absender- und Empfängerformat der Kanäle, insbesondere bei
WhatsApp durch das Präfix \texttt{whatsapp:}.

Relevante Felder sind dabei:

\begin{itemize}
    \item \texttt{From}
    \item \texttt{To}
    \item \texttt{ContentSid}
    \item \texttt{ContentVariables}
\end{itemize}

Dabei verweist \texttt{ContentSid} auf ein in Twilio hinterlegtes
Content Template. Die eigentlichen fachlichen Daten werden über
\texttt{ContentVariables} als Variablen an das Template übergeben.

Ein mögliches Template kann beispielsweise folgende Variablen verwenden:

\begin{itemize}
    \item \texttt{\{\{1\}\}} Kundinnen- oder Kundenname
    \item \texttt{\{\{2\}\}} Auftragsnummer
    \item \texttt{\{\{3\}\}} Produkt
    \item \texttt{\{\{4\}\}} Menge
    \item \texttt{\{\{5\}\}} Lieferdatum
    \item \texttt{\{\{6\}\}} Lieferzeitfenster
    \item \texttt{\{\{7\}\}} Lieferadresse
    \item \texttt{\{\{8\}\}} Bestätigungslink
\end{itemize}

Beispiel für \texttt{ContentVariables}:

\begin{lstlisting}[language={}]
{
  "1": "Max Mustermann",
  "2": "A-1001",
  "3": "Heizoel",
  "4": "3000 Liter",
  "5": "12.06.2099",
  "6": "10:00 - 11:00 Uhr",
  "7": "Beispielstrasse 12, 97070 Wuerzburg",
  "8": "https://example.com/confirmation/abc123"
}
\end{lstlisting}

Ein mögliches Nachrichten-Template kann beispielsweise wie folgt
aussehen:

\begin{lstlisting}[language={}]
Guten Tag {{1}},

für Ihren Auftrag {{2}} ist die Lieferung von {{3}} ({{4}})
am {{5}} im Zeitfenster {{6}} geplant.

Lieferadresse:
{{7}}

Bitte bestätigen Sie Ihren Liefertermin über folgenden Link:
{{8}}

Vielen Dank.
\end{lstlisting}

Beispielhaft werden daraus kanalabhängig Nachrichten mit
gleicher inhaltlicher Struktur erzeugt, wobei SMS und WhatsApp
dieselben fachlichen Werte verwenden.

% =================================================
\section{Testumgebung}
% =================================================

Für die Entwicklung und den Test der Kommunikationskanäle wird ein
Twilio-Trial-Account verwendet. Dadurch können sowohl der SMS- als auch
der WhatsApp-Versand bereits vor der Produktivsetzung getestet werden,
ohne einen produktiven Messaging-Dienst einzusetzen.

Für WhatsApp wird die \emph{Twilio WhatsApp Sandbox} genutzt. Das
jeweilige Testgerät muss zuvor in der Sandbox registriert werden, bevor
Nachrichten empfangen werden können. Der Versand erfolgt im Rahmen der
Sandbox über freigegebene Testabläufe beziehungsweise
Nachrichtentemplates.

Für den SMS-Versand wird eine von Twilio bereitgestellte Testnummer
verwendet. Im Trial-Modus können SMS ausschließlich an zuvor
verifizierte Rufnummern versendet werden. Außerdem erfolgt der Versand
nur im Rahmen des verfügbaren Trial-Guthabens.

Die erfolgreiche Zustellung kann sowohl über die Twilio-Konsole als auch
anhand der Log-Ausgaben des Backends überprüft werden.

% =================================================
\section{Produktivperspektive}
% =================================================

Für eine spätere Produktivumgebung kann derselbe technische Ansatz
beibehalten werden. Dabei müssen jedoch produktive Twilio-Ressourcen
verwendet werden, insbesondere:

\begin{itemize}
    \item produktive SMS-fähige Absendernummer
    \item freigegebener WhatsApp-Business-Sender
    \item freigegebenes Twilio-Content-Template für die genutzten Kanäle
    \item produktive Frontend-URL für Bestätigungslinks
\end{itemize}

Damit kann die Testarchitektur mit geringem Anpassungsaufwand in eine
produktive Architektur überführt werden.

% =================================================
\section{Fazit}
% =================================================

Twilio eignet sich als zentraler Messaging-Dienst für die geplante
Integration von WhatsApp und SMS. Durch die gemeinsame API, die
vergleichbare Request-Struktur und die vorhandenen Testmöglichkeiten
kann eine einheitliche technische Lösung aufgebaut werden.

Für das Projekt bietet Twilio insbesondere den Vorteil, dass sowohl
SMS als auch WhatsApp mit ähnlichen Integrationsprinzipien umgesetzt
werden können. Dadurch wird die Implementierung vereinfacht und eine
spätere Erweiterung für produktive Nutzung vorbereitet.

\end{document}
